\chapter{Giới thiệu}
\label{ch:introduction}

Chương này trình bày bối cảnh nghiên cứu, bài toán trung tâm và cơ sở thiết kế hệ thống SENTINEL---một kiến trúc nhận thức hai tầng tự động hóa quy trình phát hiện và phản hồi sự cố trong các Trung tâm Điều hành An ninh (SOC). Bằng việc phân tích hạn chế của phương pháp phòng thủ tất định truyền thống, tình trạng quá tải cảnh báo, cùng rào cản độ trễ và lỗ hổng đối kháng khi ứng dụng Mô hình Ngôn ngữ Lớn (LLM), chương này xác lập mục tiêu, câu hỏi nghiên cứu, phạm vi, phương pháp luận và các đóng góp cốt lõi của luận văn.

\section{Bối cảnh và Động lực Nghiên cứu}

Một Trung tâm Điều hành An ninh (SOC) hiện đại mỗi ngày nhận về hàng triệu sự kiện nhật ký~\cite{siem2021}, trong đó phần áp đảo là hoạt động bình thường. Vì sợ bỏ lọt, các hệ thống giám sát thường được đặt ở ngưỡng rất nhạy, nên số cảnh báo sinh ra càng nhiều và tỉ lệ báo động giả có thể vượt 80\%~\cite{alertfatigue2022}. Hệ quả là chỗ nghẽn của một SOC hiếm khi nằm ở khả năng phát hiện. Nó nằm ở khâu tiếp theo: quyết định xem cảnh báo nào đáng để một chuyên gia dừng lại xem xét. Sự chú ý của con người mới là nguồn lực khan hiếm, và nó không nhân lên được bằng cách mua thêm máy chủ.

Các công cụ tự động hoá đang dùng chưa gánh được khâu đó. Thiết bị lọc theo luật tĩnh và chữ ký cố định chỉ nhận ra những gì đã được mô tả trước, nên mọi tình huống nằm ngoài mô tả đều bị đẩy sang cho người~\cite{soar2025}. Chúng phân loại được, nhưng không giải thích được vì sao; mà việc chuyên gia phải làm lại chính là hiểu bối cảnh của một cảnh báo trước khi quyết định.

Mô hình Ngôn ngữ Lớn (LLM) hứa hẹn đúng năng lực còn thiếu ấy: đọc một sự kiện trong bối cảnh của nó và diễn giải bằng ngôn ngữ con người hiểu được. Nhưng đưa một mô hình như vậy vào giữa luồng dữ liệu vận hành thì sinh ra ba khoản chi phí mới. Thứ nhất là chi phí thời gian và tài nguyên: mỗi lần suy luận mất vài giây, trong khi dữ liệu đến liên tục~\cite{vaswani2017}. Thứ hai là chi phí niềm tin: bản thân tầng suy luận trở thành một mục tiêu, bởi nhật ký có thể mang theo câu chữ do kẻ tấn công cố ý soạn để lái mô hình~\cite{promptinjection2023}, và những gì mô hình kết luận thì phải chứng minh được là chưa bị sửa đổi khi đưa ra làm chứng cứ. Thứ ba là chi phí chủ quyền dữ liệu: nhật ký nội bộ thường không được phép rời khỏi tổ chức, nên mô hình phải chạy tại chỗ trên phần cứng hữu hạn~\cite{nist80207}.

Quan sát nền tảng dẫn tới thiết kế của luận văn là chênh lệch chi phí giữa hai loại phán quyết. Một phán quyết bằng luật thống kê tốn chưa tới một phần nghìn giây; một phán quyết bằng mô hình ngôn ngữ tốn hàng giây. Khoảng cách bốn đến năm bậc độ lớn ấy khiến việc đem mọi sự kiện qua tầng đắt nhất trở thành lãng phí có hệ thống, trong khi tuyệt đại đa số sự kiện vốn có thể kết luận được bằng những phép kiểm rất rẻ. Nghiên cứu vì thế đề xuất SENTINEL, một kiến trúc hai tầng theo nguyên lý đặt phán quyết rẻ trước phán quyết đắt: một tầng lọc nhanh giải quyết dứt điểm phần lớn lưu lượng, và một tầng suy luận chỉ được gọi tới cho phần còn lại, phía sau các rào chắn bảo vệ và một vết kiểm toán niêm phong.

Ba khoản chi phí nêu trên định hình ba hướng nghiên cứu của luận văn, và điều làm chúng trở thành bài toán nghiên cứu chứ không chỉ là bài toán kỹ thuật là ở chỗ cả ba đều phải đo được, trên dữ liệu thật, bằng những thước đo không tự tâng bốc.

\section{Mục tiêu Nghiên cứu}

Mục tiêu tổng quát của luận văn là nghiên cứu thiết kế, xây dựng nguyên mẫu và đánh giá thực nghiệm kiến trúc SENTINEL nhằm tự động hóa quy trình phát hiện, suy luận ngữ cảnh và phản hồi sự cố an ninh mạng trên hạ tầng cục bộ, đạt tối ưu về độ trễ xử lý và đảm bảo an toàn mật mã.

Mục tiêu tổng quát trên được cụ thể hóa thông qua ba nhiệm vụ nghiên cứu trọng tâm, đặt tương ứng một-một với ba câu hỏi nghiên cứu ở Mục~\ref{sec:rq}:

Thứ nhất, xây dựng đường ống lọc hai chặng phi trạng thái tại Tầng 1 vận hành ở tốc độ đường truyền, kết hợp thuật toán thống kê Welford $O(1)$, Cổng Học máy LightGBM và bộ đệm Semantic Cache Tầng 1.75 nhằm giải quyết dứt điểm phần lớn lưu lượng ngay tại tầng rẻ, không tiêu tốn tài nguyên GPU.

Thứ hai, thiết kế cơ chế bảo vệ tất định cho mô hình ngôn ngữ Tầng 2 thông qua kỹ thuật Đóng gói Dữ liệu Phân định sử dụng Nonce mật mã ngẫu nhiên, rào chắn Guardrail kiểm soát danh mục hành động đầu ra và chuỗi niêm phong HMAC-SHA256 nhằm vô hiệu hóa các cuộc tấn công tiêm nhiễm câu lệnh và bảo đảm tính toàn vẹn, phát hiện giả mạo đối với vết pháp y.

Thứ ba, phát triển cỗ máy suy luận có trạng thái dựa trên đồ thị tác tử LangGraph kết hợp Bộ nhớ Đe dọa bền vững và hệ thống RAG kép (FAISS và BM25) tích hợp tri thức MITRE ATT\&CK chuẩn STIX, cho phép sinh phán quyết có trích dẫn bằng chứng, phát hiện và vô hiệu hoá hư cấu bằng rào chắn neo-RAG tất định, đồng thời phân loại chính xác phần công việc bắt buộc phải chuyển cho chuyên gia.

\section{Câu hỏi Nghiên cứu}
\label{sec:rq}

Tiến trình thiết kế và đánh giá thực nghiệm của luận văn được dẫn dắt bởi ba câu hỏi nghiên cứu chiến lược:

Câu hỏi nghiên cứu thứ nhất (RQ1) về hiệu năng đường ống và kinh tế học tầng lọc: Làm thế nào để thiết kế một kiến trúc phân tầng có khả năng xả tải tự động luồng dữ liệu an ninh ở tốc độ đường truyền nhằm giải quyết vấn đề ``nút thắt cổ chai'' về độ trễ và tiêu thụ tài nguyên của các Mô hình Ngôn ngữ Lớn?

Câu hỏi nghiên cứu thứ hai (RQ2) về rào chắn an ninh AI và tính toàn vẹn vết pháp y: Phương pháp luận và cơ chế bảo mật nào có khả năng chống lại các rủi ro đối kháng (điển hình như tiêm nhiễm prompt qua nhật ký) nhằm bảo vệ luồng suy luận của AI, đồng thời bảo đảm tính toàn vẹn và khả năng phát hiện giả mạo của các vết chứng cứ pháp y?

Câu hỏi nghiên cứu thứ ba (RQ3) về suy luận tác tử có trạng thái và quy kết kỹ thuật: Làm thế nào để tích hợp năng lực suy luận có trạng thái và khả năng tra cứu tri thức chuyên ngành vào Tác tử AI nhằm tự động hóa quy trình phân tích, quy kết kỹ thuật tấn công và sinh ra các báo cáo pháp y minh bạch, qua đó giảm tải và phân loại chính xác phần công việc cần đến chuyên gia?

\section{Phạm vi và Đối tượng Nghiên cứu}

Đối tượng nghiên cứu trọng tâm của luận văn là Kiến trúc Nhận thức Hai tầng SENTINEL, bao gồm bộ lọc phi trạng thái Tầng 1 (thuật toán thống kê trực tuyến Welford và Cổng Học máy LightGBM), Tầng 1.75 Semantic Cache, cỗ máy tác tử có trạng thái Tầng 2 (đồ thị LangGraph và mô hình ngôn ngữ Foundation-Sec-8B-Instruct lượng tử hóa Q4\_K\_M vận hành qua llama.cpp), hệ thống RAG kép (kết hợp FAISS và BM25) tích hợp tri thức MITRE ATT\&CK, cùng các cơ chế an ninh mật mã gồm Đóng gói Dữ liệu Phân định Nonce và chuỗi niêm phong kiểm toán HMAC-SHA256.

Phạm vi đánh giá dựa trên \textbf{hai tập dữ liệu chính}: CSE-CIC-IDS2018 cho luồng giám sát mạng diện rộng và CSIC 2010 cho tấn công HTTP Tầng 7 thực tế. Hai nguồn bổ trợ---DAPT2020 cho chuỗi xâm nhập đa giai đoạn và tập zero-day---chỉ đóng vai trò \textbf{minh chứng khả thi (proof-of-concept)}, không tham gia bất kỳ tuyên bố đóng góp nào. Đánh giá đối kháng dùng 703 mẫu từ ba bộ công bố (AdvBench, deepset, jackhhao) làm tỉ lệ chính, kèm 80 mẫu tác giả biên soạn báo cáo tách bạch. Hệ thống được triển khai trên một máy trạm cục bộ (NVIDIA RTX 4060 Ti 16GB VRAM, Intel Core i7, 32GB RAM), vận hành hoàn toàn cách ly mạng. Danh mục hành động tự động giới hạn ở chặn địa chỉ IP, ghi nhật ký hoặc chuyển chuyên gia duyệt; can thiệp ngắt kết nối hạ tầng vật lý nằm ngoài phạm vi.

\section{Phương pháp Nghiên cứu}

Luận văn áp dụng phương pháp Nghiên cứu Khoa học Thiết kế (DSR) trên ba bình diện: thống kê và học máy (thuật toán Welford trực tuyến $O(1)$ tính $Z$-score kết hợp mô hình phân loại LightGBM); logic tác tử (đồ thị trạng thái LangGraph điều phối suy luận ngôn ngữ, kết hợp truy xuất lai FAISS, BM25 và hợp nhất xếp hạng RRF); và bảo mật đối kháng (thử nghiệm ứng suất chống tiêm nhiễm câu lệnh và kiểm tra toàn vẹn vết nhật ký). Hiệu quả kiến trúc được đánh giá qua Khung chỉ số 5 chiều: chất lượng phân loại, hiệu năng xử lý, khả năng kháng đối kháng, độ nhất quán của căn cứ giải thích do mô hình độc lập chấm, và toàn vẹn vết nhật ký. Ba nguyên tắc phương pháp luận được tuân thủ xuyên suốt: (i) \textbf{mẫu số phải khớp câu hỏi}---chỉ số phụ thuộc phân bố lấy từ luồng thực, chất lượng phân loại lấy từ tập cân bằng lớp; (ii) dùng thước đo miễn nhiễm lệch lớp (MCC, Balanced Accuracy) thay vì Accuracy hay $F1$ nhị phân; (iii) mọi tỉ lệ đều báo kèm mẫu số và độ phủ phép đo.

\section{Đóng góp Cốt lõi của Luận văn}

Luận văn đóng góp ba giá trị cốt lõi vào lĩnh vực an toàn thông tin và ứng dụng trí tuệ nhân tạo.

\textbf{Về kiến trúc hệ thống}, nghiên cứu đề xuất và định lượng nguyên lý \textit{đặt phán quyết rẻ trước phán quyết đắt}, đồng thời chứng minh bằng số đo rằng tỉ lệ xả tải là \textbf{hàm của tỉ lệ tấn công trong luồng}, không phải hằng số của hệ---đặc tính chưa được các công trình liên quan phát biểu tường minh.

\textbf{Về bảo mật AI}, nghiên cứu đóng góp cơ chế phòng thủ mang tính cấu trúc dựa trên Đóng gói Dữ liệu Phân định với Nonce mật mã ngẫu nhiên. Điểm cốt lõi là bảo chứng này \textbf{không phụ thuộc vào tỉ lệ chặn của bất kỳ bộ phân loại nào}: dữ liệu trong vùng bọc được xử lý như văn bản thụ động bất kể nội dung, khác hẳn các bộ lọc ngữ nghĩa xác suất vốn suy giảm trước biến thể mới.

\textbf{Về tác tử và tri thức bảo mật}, nghiên cứu thay thế kịch bản SOAR cứng nhắc bằng đồ thị tác tử có trạng thái LangGraph trên kho tri thức 433 mục MITRE ATT\&CK chuẩn STIX và Bộ nhớ Đe dọa bền vững, kèm rào chắn neo-RAG tất định buộc mọi mã kỹ thuật công bố phải hiện diện trong tài liệu truy xuất của chính lô đó---biến việc chống hư cấu từ kỳ vọng vào mô hình thành ràng buộc kiểm chứng được của kiến trúc.

\section{Cấu trúc của Luận văn}

Chương 2 hệ thống hóa cơ sở lý thuyết và tổng quan công trình liên quan nhằm xác định khoảng trống nghiên cứu. Chương 3 trình bày thiết kế kỹ thuật của SENTINEL. Chương 4 trình bày kết quả thực nghiệm định lượng, tổ chức theo đúng ba câu hỏi nghiên cứu. Chương 5 tổng kết kết quả, thừa nhận giới hạn và phác thảo hướng phát triển.
